% =============================================================================
% Introduction (unnumbered, but appears in TOC)
% -----------------------------------------------------------------------------
% The introduction in a Monash thesis-by-publication has several jobs:
%   1. State the research question and the central claims of the thesis.
%   2. Locate the work in its broader field, including a short review
%      of the relevant literature.
%   3. Motivate the choice of methods and theoretical framework.
%   4. Outline the structure of the thesis (one paragraph per chapter is
%      conventional).
%
% The introduction is typically 5,000–10,000 words. It is read by the
% examiners first and sets their expectations for the rest of the thesis.
% =============================================================================
\chapter*{Introduction}
\addcontentsline{toc}{chapter}{Introduction}
\markboth{Introduction}{Introduction}
\label{ch:introduction}

% Sample structure with placeholder section headings — feel free to depart
% from it. Most introductions have 4–8 thematic subsections like these.

\subsection*{The gap this thesis addresses}

[State the gap in current understanding that motivates the thesis. What
question does the field need answered, and why has it not yet been
answered?]

\subsection*{Why this matters}

[Outline the broader significance of the question. Who benefits from an
answer, and how?]

\subsection*{Approach taken in this thesis}

[Describe the methodological framework you adopt and why it is suited to
the question. Justify the choice over alternatives.]

\subsection*{Outline of the thesis}

This thesis is organised as follows. \Cref{ch:chapter-template} [describe
its contribution in one sentence]. The conclusion (\Cref{ch:conclusion})
[describe what the conclusion does].

% Citation example: \citet{authorYear} for in-text citations,
% \citep{authorYear} for parenthetical, \citep[p.~42]{authorYear} for
% page-specific citations. For multiple cites: \citep{a2020,b2021,c2022}.
